\hypertarget{iterators-generators-and-co-routines}{%
\chapter{ITERATORS, GENERATORS, AND
CO-ROUTINES}\label{iterators-generators-and-co-routines}}

\hypertarget{the-iterator-protocol-at-the-bytecode-level}{%
\subsection{9.1 The Iterator Protocol at the Bytecode
Level}\label{the-iterator-protocol-at-the-bytecode-level}}

An iterator is defined by implementing
\passthrough{\lstinline!\_\_iter\_\_()!} (which returns the iterator
object itself) and \passthrough{\lstinline!\_\_next\_\_()!} (which
returns the next value or raises
\passthrough{\lstinline!StopIteration!}). * \textbf{Bytecode
Compilation}: The compiler optimizes loops targeting iterators by
emitting specialized bytecodes: - \passthrough{\lstinline!GET\_ITER!}:
Pops the iterable from the stack, calls its
\passthrough{\lstinline!\_\_iter\_\_()!}, and pushes the resulting
iterator back onto the stack. - \passthrough{\lstinline!FOR\_ITER!}:
Calls \passthrough{\lstinline!\_\_next\_\_()!} on the iterator. If a
value is returned, it is pushed onto the stack and execution jumps to
the loop body. If \passthrough{\lstinline!StopIteration!} is caught, the
exception is cleared, and the instruction pointer jumps past the loop.

\hypertarget{generator-execution-frame-suspension}{%
\subsection{9.2 Generator Execution Frame
Suspension}\label{generator-execution-frame-suspension}}

Standard function calls allocate a stack frame, execute, and release the
frame. Generator functions behave differently: * \textbf{Heap
Preservation}: When a function containing the
\passthrough{\lstinline!yield!} keyword is called, it returns a
generator object wrapping a \passthrough{\lstinline!PyFrameObject!}. *
\textbf{Yield Mechanism}: When \passthrough{\lstinline!yield!} is
evaluated, the interpreter saves the current state (evaluation stack
values and the next instruction pointer
\passthrough{\lstinline!f\_lasti!} in the frame) and returns the yielded
value. The frame is \textbf{not} popped from the call stack; it remains
allocated on the heap. * \textbf{Next Mechanism}: When
\passthrough{\lstinline!next()!} is called again on the generator, the
CPython evaluation loop loads the suspended frame, restores its
registers and evaluation stack, and continues execution from the saved
\passthrough{\lstinline!f\_lasti!}.

\begin{lstlisting}
Generator Frame Suspension Flow:
[Call generator] -> [Allocates PyFrameObject on Heap]
                           |
                           v
[Run bytecode] ------> [Hits YIELD]
   ^                       |
   |                       v
   |                  [Save frame state (f_lasti)] -> [Return value to caller]
   |                       |
   +---[Call next()]-------+
\end{lstlisting}

\hypertarget{generator-communication-protocol}{%
\subsection{9.3 Generator Communication
Protocol}\label{generator-communication-protocol}}

Generators support bidirectional data exchange via the methods: *
\passthrough{\lstinline!send(value)!}: Resumes execution and passes
\passthrough{\lstinline!value!} back into the generator as the result of
the \passthrough{\lstinline!yield!} expression. *
\passthrough{\lstinline!throw(type, value)!}: Raises an exception at the
generator's current execution point. The generator can catch this
exception or propagate it. * \passthrough{\lstinline!close()!}: Raises a
\passthrough{\lstinline!GeneratorExit!} exception in the generator to
run cleanup blocks (e.g.~\passthrough{\lstinline!try/finally!} resource
releases) and terminates the generator.

\begin{lstlisting}[language=Python]
# Bidirectional Communication Trace
def bidirectional_generator():
    try:
        val = yield "Ready"
        print("Received value inside generator:", val)
        yield f"Echo: {val}"
    except ValueError:
        print("Caught ValueError in generator")
        yield "Recovered"

gen = bidirectional_generator()
print(next(gen)) # Start generator (yields "Ready")
print(gen.send("Hello World")) # Sends "Hello World", yields "Echo: Hello World"

gen_err = bidirectional_generator()
next(gen_err)
print(gen_err.throw(ValueError)) # Raises error inside generator, yields "Recovered"
\end{lstlisting}

\hypertarget{pep-380-generator-delegation-via-yield-from}{%
\subsection{9.4 PEP 380: Generator Delegation via yield
from}\label{pep-380-generator-delegation-via-yield-from}}

Introduced in Python 3.3,
\passthrough{\lstinline!yield from <iterable>!} delegates operations to
another generator (the sub-generator): * \textbf{Transparent Channel}:
It establishes a direct, bidirectional channel between the outer caller
and the sub-generator. * \textbf{Automated Propagation}: - Values sent
via \passthrough{\lstinline!send()!} are routed directly to the
sub-generator. - Exceptions thrown via \passthrough{\lstinline!throw()!}
are raised inside the sub-generator. -
\passthrough{\lstinline!StopIteration!} raised by the sub-generator is
caught automatically, and its \passthrough{\lstinline!value!} attribute
becomes the result of the \passthrough{\lstinline!yield from!}
expression.

\begin{center}\rule{0.5\linewidth}{0.5pt}\end{center}
